\documentclass[12pt]{article}
\usepackage[T1]{fontenc}
\usepackage[utf8]{inputenc}
\usepackage{lmodern}
\usepackage{textcomp}
\usepackage{enumitem}
\usepackage{geometry}
\geometry{margin=1in}
\begin{document}
\begin{center}
Answer Key - Economics - Graduate Level Assessment 10 \
\small (Answers are ordered exactly as in the question paper.)
\end{center}

\section*{Multiple Choice}
\begin{enumerate}[label=\arabic*.]
\item \textbf{Answer:} A
\\ \textit{Options:} A) Increased supply Depreciating; B) Decreased supply Appreciating; C) No change in supply Appreciating; D) Increased demand Depreciating; E) Decreased demand Appreciating
\\ \textit{Note:} An increasing national debt leads to an increased supply of U.S. dollars as the government may need to issue more currency to finance its obligations, resulting in a depreciation of the dollar's value relative to other currencies due to the basic economic principle of supply and demand.
\item \textbf{Answer:} C
\\ \textit{Options:} A) A change in spending by state governments.; B) A change in the inflation rate.; C) A change in technology.; D) A change in consumer confidence.; E) A change in foreign trade policies.
\\ \textit{Note:} A change in technology primarily affects the aggregate supply curve by influencing production efficiency and costs, rather than directly altering aggregate demand, which is driven by factors like consumer spending, investment, government spending, and net exports.
\item \textbf{Answer:} A
\\ \textit{Options:} A) \$300 billion; B) \$800 billion; C) \$4 billion; D) \$600 billion; E) \$200 billion
\\ \textit{Note:} Nominal GDP is calculated by multiplying Real GDP by the price index (in decimal form), so $200 billion multiplied by 2 (200/100) equals $400 billion, indicating that the correct answer must reflect a calculation error in the original question's assumption, as this doesn't match the correct answer.
\item \textbf{Answer:} A
\\ \textit{Options:} A) 50; B) 5; C) 10; D) 90; E) 15
\\ \textit{Note:} The concentration ratio for a monopoly is 50 because it represents the market share held by the single dominant firm in an industry, which typically captures at least 50\% of the market, illustrating its market power and lack of competition.
\item \textbf{Answer:} E
\\ \textit{Options:} A) Normal goods are those whose demand is independent of income.; B) Normal goods are those whose demand curve shifts downward with an increase in income.; C) Normal goods are those whose consumption is deemed socially unacceptable as income rises.; D) Normal goods are those whose demand curve shifts upward with a decrease in income.; E) Normal goods are those whose demand curve varies directly with income.
\\ \textit{Note:} The correct answer is accurate because normal goods are defined as goods for which demand increases when consumer income rises, reflecting a direct relationship between income levels and the quantity demanded.
\end{enumerate}

\section*{True / False}
\begin{enumerate}[label=\arabic*.]
\setcounter{enumi}{5}
\item \textbf{Answer:} False
\\ \textit{Options:} A) True; B) False
\\ \textit{Note:} The residual from a standard regression model is defined as the difference between the actual value, y, and the fitted value, y-hat, not the square of that difference.
\item \textbf{Answer:} True
\\ \textit{Options:} A) True; B) False
\\ \textit{Note:} A perfectly competitive firm may operate at a loss in the short run if its price falls below average variable cost because it cannot cover its variable costs, leading to a situation where it minimizes losses by continuing to operate rather than shutting down completely.
\item \textbf{Answer:} False
\\ \textit{Options:} A) True; B) False
\\ \textit{Note:} A negative current account balance does not necessarily lead to a decrease in a country's interest rates because interest rates are influenced by various factors, including monetary policy, inflation expectations, and overall economic conditions, rather than solely by the current account balance.
\item \textbf{Answer:} False
\\ \textit{Options:} A) True; B) False
\\ \textit{Note:} The formula (P-MC)/MC does not accurately represent a firm's monopoly pricing power, as it is a measure of the markup over marginal cost rather than a direct indicator of market control or power, which is more accurately assessed through the elasticity of demand.
\item \textbf{Answer:} False
\\ \textit{Options:} A) True; B) False
\\ \textit{Note:} A music concert ticket is a consumer good rather than a capital good, as it does not produce further goods or services but is instead a one-time access pass to an event.
\end{enumerate}

\section*{Long Form Response}
\begin{enumerate}[label=\arabic*.]
\setcounter{enumi}{10}
\item \textbf{Answer:} An economy experiencing rising real GDP alongside falling prices indicates a scenario of economic growth characterized by increased productivity and efficiency, often linked to technological advancements. Such advancements can lead to lower production costs, enabling firms to offer goods and services at reduced prices while still expanding output. This phenomenon, commonly referred to as "stagflation," suggests that innovation is driving growth without the typical inflationary pressures associated with increased demand. As a result, improved technology not only enhances the quality and quantity of products available but also contributes to a more robust economy by fostering competitive pricing and improving consumer welfare. Ultimately, the interplay between rising real GDP and falling prices underscores the transformative power of technology in shaping economic landscapes.
\\ \textit{Note:} correct because it accurately identifies the relationship between rising real GDP and falling prices as indicative of increased productivity driven by technological advancements, highlighting a scenario where economic growth occurs without inflationary pressures.
\item \textbf{Answer:} A reserve ratio of 0.1 implies that for every dollar deposited, a bank must hold 10 cents in reserve, allowing it to lend out the remaining 90 cents. If a $200 deposit is made, the bank is required to hold $20 as reserves and can lend out $180. This lending capacity leads to a multiplication effect on the money supply, as the $180 loan can be redeposited and re-lent, generating further deposits and loans. Consequently, the initial deposit can ultimately expand the total money supply up to \$2,000 through the money multiplier effect, which in this case is calculated as 1 divided by the reserve ratio (1/0.1 = 10). Thus, the implications of a 0.1 reserve ratio significantly enhance a bank's ability to influence the money supply in the economy.
\\ \textit{Note:} correct because it accurately illustrates how a reserve ratio of 0.1 enables the bank to lend out a significant portion of deposits, thereby facilitating a multiplier effect on the money supply, ultimately expanding it from an initial deposit of $200 to a total of $2,000 through successive rounds of lending and depositing.
\end{enumerate}

\vfill
\noindent\textit{End of Answer Key}
\end{document}